% introduction.tex (Paper II)
% Frames the central claim (Eq. (137) of Paper I) and includes the
% audit table (Table \ref{tab:audit}) so the reader sees the
% PASS/STUB decomposition before any dependent claim.

\label{sec:introduction}

Paper~I (\emph{Geometry First}) established the bipartite octahedral
causal lattice $\mathcal{T}_\diamond^3$ with the unity constraint
$\mathcal{A}=1$ as a discrete substrate from which the Dirac equation,
the Born rule, hydrogen quantization at four significant figures,
and the structural form of the induced gauge action are derived.
The paper closes (§15) with a precise central conjecture, repeated
here as the anchor of Paper~II:
%
\begin{equation}
\label{eq:automorphism_conjecture}
\mathrm{Aut}(\mathcal{T}_\diamond^3, \mathcal{A}=1)
\;=\; SO(3,1) \times SU(3) \times SU(2) \times U(1).
\end{equation}
%
Both sides are precise finite-dimensional Lie algebras with real
dimension $18 = 6 + 8 + 3 + 1$.  The match is more than dimensional:
the structure constants of the lattice's automorphism algebra must
reproduce those of the right-hand side factor by factor.  Either the
automorphism calculation closes the equality (and the Standard Model
gauge group is derived from a single conservation axiom on a discrete
substrate), or it does not (and Paper~II is a precise statement of
the obstruction).  Both outcomes are concrete results.

\subsection{Status of the conjecture entering Paper~II}
\label{subsec:status_at_entry}

Paper~I's v1.0 release supplies six computational scaffolding
scripts in \nolinkurl{src/utilities/}
(\nolinkurl{automorphism_*.py}, \nolinkurl{tick_rule_*.py}) that
establish the load-bearing
sub-claims of Eq.~\eqref{eq:automorphism_conjecture}.  On the
existing per-site amplitude $\mathbb{C}^2 = (\psi_R, \psi_L)$ the
direct product $SO(3,1) \times U(1)$ (dim 7) is verified, and the
proposed per-site $SU(2)$ generators on this carrier are confirmed
to coincide with the Lorentz rotation subgroup --- so
$SU(2)_W$ and $SU(3)$ require structural extensions of the per-site
amplitude.  On the extended amplitude
$\mathbb{C}^{12} = \mathbb{C}^2 \otimes \mathbb{C}^2 \otimes
\mathbb{C}^3$ the four conjecture factors commute pairwise (dim 18
verified), the bipartite tick rule extends consistently
($\mathcal{A}=1$, parity, and $SU(2)_W \times SU(3)$ all preserved),
and the bipartite Wilson plaquette $V_1, -V_2, -V_1, V_2$ is
gauge-invariant.  The audit table (Table~\ref{tab:audit}) records
each sub-claim with its supporting script.

%% The audit table is included immediately after the status section
%% so the reader meets it before any claim that depends on a status
%% entry.
% ============================================================
% audit_table.tex (Paper II)
%
% Maps Paper II's claims to the supporting derivations and to the
% sympy verification scripts in src/utilities/.
%
% The central claim is the conjecture stated in Paper~I as Eq.~(137):
%   Aut(T_diamond, A=1) = SO(3,1) x SU(3) x SU(2) x U(1)
% (real dim 18 = 6+8+3+1).  Paper II's task is to close this row by
% showing the four established sub-claims plus the three open
% sub-claims combine to the equality.
%
% Status semantics (paste into the abstract or front-matter as well):
%   PASS  -- derivation complete or verification confirms the claim.
%   PART  -- mechanism demonstrated, full quantitative match pending.
%   STUB  -- analytical result with no numerical confirmation yet,
%            or planned calculation not yet done.
%   FAIL  -- tested and disconfirmed.
%
% Use \texttt{} for code/script names and file paths inside cells.
% ============================================================

\label{tab:audit}

{\scriptsize
\setlength{\tabcolsep}{4pt}
\renewcommand{\arraystretch}{0.95}
\begin{longtable}{@{}p{3.2cm}p{3.6cm}p{3.4cm}p{3.8cm}p{1.0cm}@{}}
\caption{\small System audit: Paper~II claims mapped to supporting
derivations and sympy verification scripts.  Status semantics
(\texttt{PASS} / \texttt{PART} / \texttt{STUB} / \texttt{FAIL}) per
the front-matter.  The central claim is the bottom row;
its status is the conjunction of the rows above.}
\label{tab:audit_caption} \\
\hline
\textbf{Claim} & \textbf{Mechanism} & \textbf{Continuum / formal limit} & \textbf{Evidence} & \textbf{Status} \\
\hline
\endfirsthead

\multicolumn{5}{l}{\textit{(Table~\ref{tab:audit}, continued from previous page)}} \\
\hline
\textbf{Claim} & \textbf{Mechanism} & \textbf{Continuum / formal limit} & \textbf{Evidence} & \textbf{Status} \\
\hline
\endhead

\hline
\multicolumn{5}{r}{\textit{(continued on next page)}} \\
\endfoot

\hline
\endlastfoot

% ── Established factors ─────────────────────────────────────────────
Discrete spatial $\mathrm{Aut}(\Gamma, \mathbf{V})$ has order 48
    & Enumeration of basis-preserving linear maps on
      $\{\pm\mathbf{V}_1, \pm\mathbf{V}_2, \pm\mathbf{V}_3\}$
    & Cubic point group $B_3 \cong O_h$;
      12 elements orthogonal in $\mathbb{R}^3$
    & \nolinkurl{automorphism_discrete.py} (sympy)
    & \texttt{PASS} \\

RGB sublattice symmetry contributes only $\mathbb{Z}_3 \subset SU(3)$
    & Cyclic shift of the three RGB basis vectors generates a
      $\mathbb{Z}_3$ subgroup
    & Abelian; cannot generate non-abelian $SU(3)$ on its own
    & \nolinkurl{automorphism_rgb_su3.py} (sympy)
    & \texttt{PASS} \\

$SO(3,1) \times U(1)$ acts on the existing $\mathbb{C}^2 = (\psi_R, \psi_L)$
    & Lorentz from $O_h$-averaged Dirac (Paper~I §6); $U(1)$
      from the per-site phase oscillator
    & Real dim $6 + 1 = 7$, direct product on the existing
      per-site amplitude
    & \nolinkurl{automorphism_direct_product.py} (sympy);
      verifies $[\mathfrak{so}(3,1), \mathfrak{u}(1)] = 0$ and that
      the proposed per-site $SU(2)$ generators are the same matrices
      as the Lorentz rotation subgroup
    & \texttt{PASS} \\

Direct-product structure on extended $\mathbb{C}^{12}$
    & Per-site $\mathbb{C}^2 \otimes \mathbb{C}^2 \otimes \mathbb{C}^3$
      with the four factors acting on disjoint tensor slots
    & Real dim $6 + 8 + 3 + 1 = 18$;
      $\mathfrak{so}(3,1) \oplus \mathfrak{su}(3) \oplus
       \mathfrak{su}(2) \oplus \mathfrak{u}(1)$ commutes pairwise
    & \nolinkurl{automorphism_direct_product_extended.py} (sympy)
    & \texttt{PASS} \\

Tick-rule consistency on $\mathbb{C}^{12}$
    & Trivial tensor extension
      $T_\text{ext} = T_\text{chirality} \otimes I_2 \otimes I_3$
      of the bipartite tick rule
    & Preserves $\mathcal{A}=1$, bipartite parity, and the global
      $SU(2)_W \times SU(3)$ symmetries
    & \nolinkurl{tick_rule_extended_consistency.py} (sympy);
      four checks pass
    & \texttt{PASS} \\

Wilson plaquette gauge invariance
    & Bipartite Wilson plaquette
      $V_1, -V_2, -V_1, V_2$ (smallest non-trivial closed loop;
      no triangular plaquettes since $V_1 + V_2 + V_3 \neq 0$)
    & $\mathrm{Tr}(W) = 2\cos(a_1 + a_2 + a_3 + a_4)$
      gauge-invariant by cyclic property
    & \nolinkurl{tick_rule_gauge_invariance.py} (sympy);
      $U(1)$ explicit, $SU(2)$ sample, general non-abelian by
      $V V^\dagger = I$
    & \texttt{PASS} \\

SU(3) representation branch consistency
    & Only the $\mathbf{3} \oplus \mathbf{3}$ interpretation
      (Branch~A: $T_a = \lambda_a / 2$ on every site) commutes
      with the bipartite tick rule
      $T_\text{chir} \otimes I_2 \otimes I_3$; the
      $\mathbf{3} \oplus \bar{\mathbf{3}}$ interpretation
      (Branch~B: RGB in $\mathbf{3}$, CMY in $\bar{\mathbf{3}}$)
      obstructs on the 5 real Gell-Manns
      $\{\lambda_1, \lambda_3, \lambda_4, \lambda_6, \lambda_8\}$
    & Residual factors as $[T_\text{chir}, P_R] \otimes I_2 \otimes
      \mathrm{Re}(\lambda_a)$; forces the bipartite parity to be
      linear under the existing tick rule (antilinear $CP$
      candidates require modifying the tick rule itself)
    & \nolinkurl{su3_branch_consistency.py} (sympy);
      8/8 Branch~A commutators vanish, 3/8 Branch~B commute,
      5/8 Branch~B obstruct
    & \texttt{PASS} \\

Non-abelian $SU(3)$ from the $\mathbb{C}^3$ colour memory
    & The colour-memory tick rule applies $U_j$ on $\mathbb{C}^3$ on
      a tick along $\mathbf{V}_j$; the $S_3$ orbit $\{X_j = -i\log U_j\}$
      closes under Lie brackets to a subalgebra of $\mathfrak{su}(3)$.
      The natural ``rotate toward $|j\rangle$'' reading
      ($X_1 = \lambda_1 + \lambda_4$ for $j=1$) and any
      mixed Cartan-plus-real-off-diagonal candidate generate the
      full $\mathfrak{su}(3)$
    & Degenerate cases (purely diagonal Cartan only, or purely
      imaginary antisymmetric off-diagonals) yield strict
      subalgebras: $\mathfrak{h}$ (Cartan, dim 2) or
      $\mathfrak{su}(2)$ (dim 3) respectively.
      The discrete RGB $\mathbb{Z}_3$ of
      \nolinkurl{automorphism_rgb_su3.py} is recovered as
      the abelian-Cartan corner.  The $\mathfrak{su}(3)$ factor of
      Eq.~(137) is dynamically generated, not merely admitted
    & \nolinkurl{su3_generation_from_colour_memory.py} (sympy);
      8 candidate $X_1$'s surveyed, closure dimensions
      $\{0, 2, 2, 3, 3, 3, 8, 8\}$
    & \texttt{PASS} \\

\hline
% ── Open sub-claims ─────────────────────────────────────────────────
Exact equality vs containment in $\mathrm{Aut}_\text{ext}$
    & The discrete-Hermitian centralizer of the bipartite tick
      rule on $\mathbb{C}^{12}$ (= operators $X$ commuting with
      $\sigma_x \otimes I_2 \otimes I_3$) has dimension 71 in
      $\mathfrak{su}(12)$, equivalently $\mathfrak{u}(6) \oplus
      \mathfrak{u}(6) - 1$.  Of the 14 Hermitian generators of
      Eq.~(137), exactly 12 ($J_1$, 3 $SU(2)_W$, 8 $SU(3)_c$)
      lie in the centralizer; the remaining 2 ($J_2, J_3$ of the
      Lorentz rotations) are continuum-emergent
    & 59 extra generators commute with the tick rule but are
      \emph{factor-mixing} operators -- couplings like
      $I_2 \otimes \sigma_a \otimes \lambda_b$ (isospin-colour mixing,
      24 generators), $\sigma_x \otimes \sigma_a \otimes I_3$
      (3), $\sigma_x \otimes I_2 \otimes \lambda_a$ (8), and
      $\sigma_x \otimes \sigma_a \otimes \lambda_b$ (24).
      These violate the SM's factor-product gauge invariance.
      Equality $=$ in Eq.~(137) holds iff the factor-product
      structure is imposed as an additional constraint
    & \nolinkurl{automorphism_centralizer_extended.py} (sympy);
      143 $\mathfrak{su}(12)$ generators enumerated, exactly 71
      commute, 12 are conjecture-aligned and 59 are factor-mixings
    & \texttt{PART} \\

SM-chirality coupling alignment
    & Under Branch~A SU(3) and the existing tick rule, the unique
      linear bipartite parity is $P = \sigma_x \otimes I_2 \otimes I_3$
      (by Schur on the irreducible $\mathbf{3}$).
      $P$ anticommutes with $\gamma_5 = \sigma_z \otimes I_2 \otimes I_3$
      and swaps the SM's $\psi_L \leftrightarrow \psi_R$ (spatial parity,
      not chirality projection);
      $P$ commutes with the proposed $SU(2)_W$ generators (vector-like,
      not the SM's chiral coupling).
      Bipartite $\mathbb{Z}_2$ and SM chirality $\mathbb{Z}_2$ are
      orthogonal Bloch involutions on the chirality $\mathbb{C}^2$
    & SM left-handed doublets / right-handed singlets are
      empirically required; framework provides instead a vector-like
      $SU(2)$ coupling and a spatial-parity $\mathbb{Z}_2$ that is
      orthogonal to (not equal to) the SM's chirality $\mathbb{Z}_2$.
      The follow-on question (whether a modified antilinear tick
      rule admits SM chirality / $CP$) is settled negatively below
    & \nolinkurl{chirality_parity_alignment.py} (sympy);
      5 candidates tested, unique survivor $M = I_3$ (and its
      projective equivalent $-I_3$); all 8 structural identities
      verified
    & \texttt{PART} \\

Modified tick rule for Branch~B SU(3) compatibility
    & Ansatz $T_\text{ext}^B = is \cdot I_{12} + c \cdot (\sigma_x
      \otimes I_2 \otimes C) \circ K$ with $K$ global complex
      conjugation: requires the colour matrix $C$ to anticommute with
      every Gell-Mann generator $\lambda_a$
    & No non-zero $C$ exists: anticommutation with $\lambda_3$
      restricts $C$ to entries $C_{12}, C_{21}, C_{33}$;
      anticommutation with $\lambda_8$ then forces those to zero.
      Representation-theoretically: $\mathbf{3}$ and $\bar{\mathbf{3}}$
      of SU(3) are inequivalent irreps; the unique antilinear
      intertwiner is global colour conjugation, which does not
      satisfy the chirality-dependent commutativity Branch~B
      demands.  Lattice cannot incorporate SM-style $CP$ as a
      discrete symmetry of any natural tick modification
    & \nolinkurl{tick_rule_cp_modified.py} (sympy);
      9 candidates tested, all fail; structural proof via
      $\lambda_3, \lambda_8$ anticommutation;
      isospin entanglement does not relax the colour constraint
    & \texttt{FAIL} \\

Explicit $1/g^2$ prefactor for $SU(2)_W$, $SU(3)$ Wilson actions
    & Non-abelian generalisation of Paper~I's bipartite-plaquette
      calculation: $1 - \tfrac{1}{N}\mathrm{Re}\,\mathrm{Tr}\,W_{ab} =
      \tfrac{a^4 T_F}{2N}(V_a^i V_b^j F^c_{ij})^2$ with $T_F = 1/2$.
      The Q-tensor (eigenvalues $\{4, 4, 16\}$) and the
      $O_h$-averaged Maxwell-style density are inherited from
      Paper~I unchanged
    & Spectator-factor counting on the per-site
      $\mathbb{C}^2 \otimes \mathbb{C}^2 \otimes \mathbb{C}^3$
      gives $N_f^{SU(2)} = 6$, $N_f^{SU(3)} = 4$; sharp lattice-
      scale prediction $g_3^2/g_2^2 = 3/2$ independent of the
      universal one-loop prefactor $c$ Paper~I left open.
      Universal $c$ remains open (the explicit
      $-\mathrm{Tr}\ln D_\text{lat}[U]$ calculation is a
      paper-length follow-on for both papers)
    & \nolinkurl{induced_gauge_action_nonabelian.py} (sympy);
      Q-tensor verified, spectator counts derived, ratio
      $g_1^2 : g_2^2 : g_3^2 = 1 : 4 : 6$ at lattice scale
    & \texttt{PART} \\

\hline
% ── Central claim ───────────────────────────────────────────────────
\textbf{Eq.~(137) of Paper~I:}
$\mathrm{Aut}(\mathcal{T}_\diamond^3, \mathcal{A}=1) = SO(3,1) \times SU(3) \times SU(2) \times U(1)$
    & Conjunction of the four established factors above plus the
      three open sub-claims; structure constants must match factor
      by factor, not just dimensions
    & Standard Model gauge group plus Lorentz, derived from a
      single conservation axiom on a discrete substrate
    & Resolves favourably iff the three open sub-claims resolve
      favourably; characterises the obstruction otherwise
    & \texttt{STUB} \\

\hline
\end{longtable}
}


\subsection{Open subproblems}
\label{subsec:open_subproblems}

What Paper~II takes up:

\begin{enumerate}
\item \textbf{Exact equality vs containment} in
$\mathrm{Aut}_\text{ext}$.  The verified statement is
$\supseteq$; the conjecture asserts $=$.  Whether the extended
structure admits additional symmetries beyond the four-factor
product is a finite Lie-algebra enumeration question constrained by
commutation with the established factors.

\item \textbf{SM-chirality coupling alignment.}  The Standard Model's
$SU(2)_W$ couples only to left-handed fermion doublets; right-handed
fermions are $SU(2)_W$ singlets.  Whether the framework's bipartite
RGB/CMY parity is the SM's chirality projector, and whether the
proposed $SU(2)_W$ on the per-site weak-isospin $\mathbb{C}^2$
couples asymmetrically to $\psi_R$ versus $\psi_L$ as required by
the SM, is the load-bearing structural-alignment problem.

\item \textbf{Explicit $1/g^2$ prefactor} for the $SU(2)_W$ and
$SU(3)$ Wilson actions on the bipartite octahedral lattice.  The
existing $U(1)$ induced-gauge-action calculation (Paper~I
Appendix~B) supplies the template; the non-abelian generalisation
goes through with matrix-valued link variables.
\end{enumerate}

\subsection{Scope and outcomes}
\label{subsec:scope}

Paper~II's scope is precisely the three subproblems above.  All
other open questions noted in Paper~I --- the Farey-sequence mass
spectrum (Paper~I follow-on \#3), the Veneziano connection
(\#3), the proton-interior structure (\#16 dependency), plasma /
gradient ionization (\#1), and the like --- are out of scope and
deferred to subsequent papers in the series.

The paper has two viable terminal states.  If the three
subproblems resolve in the same direction, the conjecture closes and
the Standard Model gauge group is derived; the paper publishes that
derivation.  If any subproblem fails, the paper publishes a precise
statement of the obstruction and the framework's verified scope
(\emph{relativistic Hilbert-space + gravity + Born-rule structures}
plus $SO(3,1) \times U(1)$ on the existing per-site $\mathbb{C}^2$,
per Paper~I) is restated unchanged.
